\section{Relativistic-Driven Dynamics}
\label{sec:relativistic_dynamics}

The DRM dynamics used as theoretical inspiration introduces a nonlinear flow
controlled by a relativistic scale parameter \cite{muniz_relativistic_drm}. In
this formulation, a local velocity norm $v$ and a limit constant $c$ define a
Lorentz-like factor
\begin{equation}
  \gamma(v) = \frac{1}{\sqrt{1-\|v\|^2/c^2}}.
\end{equation}
As $\|v\|$ approaches $c$, the scale factor grows, amplifying the geometric
effect of fast transitions while keeping the system bounded by the velocity
limit.

In the DRM Transformer, we do not simulate a physical spacetime. The Lorentz
factor is used as a differentiable modulation of attention distances. Let
$q_t$ and $k_s$ be manifold coordinates of a query and a key. We estimate a
local speed from coordinate differences and compute
\begin{equation}
  \tilde d(q_t,k_s) = \gamma^\alpha d_G(q_t,k_s),
\end{equation}
where $\alpha$ controls the strength of gamma-scaling. The score used in
attention is proportional to $-\tilde d^2$.

This mechanism has two intended effects. First, it makes fast semantic
transitions more expensive, acting as a geometric stabilizer. Second, it offers
a controlled way to introduce nonlinear scaling into the attention field
without changing the autoregressive interface of the model.
